\section{考察}

\subsection{サムネイル選択手法について}

\subsubsection{採用した代表フレーム選択基準の特徴}

課題 2 では、中央フレームを選ぶ方法ではなく、動画全体の平均ヒストグラムに最も近い
フレームを選ぶ方法を採用した。この基準の利点は、選択結果が「動画のどの位置か」では
なく「動画の内容」に基づくことである。中央フレーム法は実装が単純だが、中央が
たまたま場面転換・手ぶれ・フェードの途中であった場合にも、その品質を評価せずに
採用してしまう。一方、平均ヒストグラム法では、例外的な明るさ分布を持つフレームは
平均から遠ざかるため選ばれにくく、動画の典型的な画面構成を持つフレームが選ばれる。

実際、People - 84973.mp4(総フレーム数 241)では中央フレームは 120 であるが、
採用手法はフレーム 185 を選択した。すなわち両手法の選択は一致せず、
採用手法は動画後半の人物が大きく写る場面を「より平均的な画面」と判定した。

\subsubsection{採用手法の限界}

平均ヒストグラム法はグレースケールの明るさ分布のみを比較するため、
構図や被写体の意味内容は考慮しない。明るさ分布が平均的でも構図として
無意味なフレームが選ばれる可能性は残る。また、複数の場面からなる動画では
「全場面の平均」がどの場面とも異なる分布になり、代表性が損なわれることが
考えられる。この点は、場面ごとに代表フレームを選ぶ(課題 7 の変わり目検出と
組み合わせる)ことで改善できると考えられる。

\subsection{スペクトログラムのパラメータについて}

\subsubsection{メッセージが読みやすかった設定}

窓長({\tt nperseg})1024、オーバーラップ({\tt noverlap})768(窓長の 75\%)、
ハン窓、dB 表示の組み合わせで、メッセージ「Good work」が最も明瞭に判読できた。
標本化周波数 44100 Hz、窓長 1024 のとき、周波数分解能は
$44100 / 1024 \approx 43$ Hz、時間分解能(ホップ幅 256 サンプル)は約 5.8 ms であり、
約 6 秒の音声に対して横方向約 1000 列の十分な密度が得られる。

\subsubsection{パラメータが不適切な場合との違い}

窓長を極端に小さくすると周波数分解能が粗くなり、文字の縦方向の形状がつぶれる。
逆に窓長を極端に大きくすると時間分解能が粗くなり、文字の横方向の形状がつぶれる。
これは式 \ref{eq:stft} の窓長 $N$ に関する時間分解能と周波数分解能の
トレードオフそのものである。また、dB 変換(式 \ref{eq:db})を行わず線形スケールで
表示すると、パワーの大きい少数の成分だけが濃く表示され、メッセージ部分との
コントラストが失われて判読できなくなる。

\subsection{n-gram と stop words について}

\subsubsection{stop words 除去の効果}

図 \ref{fig:2gram} と図 \ref{fig:2gram_sw} の比較から、除去前は上位 10 個のうち
「in the」「of the」「as a」「with the」など機能語を含む組が半数を占めるのに対し、
除去後は「natural language」「language processing」「machine translation」
「artificial intelligence」など、テキストの主題(自然言語処理)を直接表す
内容語の組だけが残った。このことから、stop words 除去は 2-gram 分析において
文章の内容的特徴を浮かび上がらせる効果があるといえる。

\subsubsection{トークン化の設計と「e g」の出現}

除去後の上位に「e g」という 2-gram が現れた(9 回)。これは原文中の略語
「e.g.」がトークン化(正規表現 {\tt [a-z]+})によって「e」と「g」に分割された
ものである。1 文字語「e」「g」は設計した stop words 集合に含めなかったため、
両者が隣接する 2-gram として集計された。内容分析の目的からは「e.g.」は
機能語に近く、stop words に 1 文字語や略語を追加するか、トークン化の段階で
「e.g.」を 1 語として扱う設計が望ましいと考えられる。この点は stop words の
設計がトークン化の仕様と不可分であることを示している。

\subsubsection{2-gram 構成方式の影響}

本実装では stop words を除去した後の単語列から 2-gram を構成した。このため、
原文では stop words を挟んで離れていた 2 語(例:「generation of natural」の
「generation」と「natural」)が新たな隣接ペア「generation natural」として
集計される。実際、図 \ref{fig:2gram_sw} には原文に連続表現として存在しない
「generation natural」が出現している。共起関係を広く拾う目的にはこの方式が
適するが、原文に実在する連接表現のみを数えたい場合は、2-gram を構成してから
stop words を含む組を捨てる方式を採るべきである。両方式は目的に応じて
使い分けるべき設計選択である。

\subsection{変わり目検出の閾値決定法について}

\subsubsection{3 方式の比較}

表 \ref{tab:detect} のとおり、平均 $+3\sigma$ 方式・上位 3 ピーク方式・
移動平均基準方式の検出結果は完全に一致した。これは図 \ref{fig:framediff} に
見られるように、変わり目の変化量(最大で約 $2.3 \times 10^8$)がショット内の
変化量(概ね $10^7$ 以下)より 1 桁以上大きく、どの合理的な閾値でも分離できる
ためである。

ただし 3 方式の性質は異なる。上位 $n$ ピーク方式は変わり目の個数が既知で
なければ使えず、汎用性に欠ける。平均 $+k\sigma$ 方式は個数の事前知識を
要しないが、動画全体で 1 つの閾値を使うため、動きの激しいショットと静かな
ショットが混在すると誤検出・見逃しが起きやすい。移動平均基準方式は閾値が
局所的な変化量に追従するため、この問題に最も頑健である。本課題の動画では
差が現れなかったが、ショット内の動きが大きい実映像では方式間の差が現れると
考えられる。

\subsubsection{検出結果の妥当性}

図 \ref{fig:boundary} の目視確認により、検出された 3 フレームすべてで場面が
完全に入れ替わっていることを確認した。誤検出(変わり目でないフレームの検出)も
見逃し(4 区間構成に必要な 3 境界のうち検出されなかったもの)もなく、
本課題の動画に対しては単純なフレーム間差分の総和で十分な精度が得られた。

\subsection{関数ベース設計とクラスベース設計の比較}

\subsubsection{データと処理の凝集}

課題 1〜7 のような関数ベースの設計では、動画のパス・タイトル・ID・サムネイルと
いった関連するデータを関数間で個別に受け渡す必要があり、扱う動画が増えると
変数管理が煩雑になる。課題 8 の Video クラスでは、これらが 1 つのインスタンスに
凝集され、「その動画に対する操作」がメソッドとして明示されるため、
複数の動画を扱うプログラムでも見通しがよい。{\tt get\_info()} のように
状態を辞書としてまとめて返す窓口を用意することで、クラス内部の変数名の変更が
利用側に波及しにくくなる利点もある。

\subsubsection{高階関数による処理の分離}

色反転動画の生成を「フレーム変換」と「動画入出力」に分離したことで、
二値化・グレースケール化の追加はフレーム変換関数を 1 つ書くだけで済んだ。
仮に分離しない設計であれば、動画の読み込み・リサイズ・書き出しという同一の
コードを処理の種類だけ複製することになり、コーデック変更などの修正が
全コピーに波及する。処理の差し替え可能性という点で、高階関数方式は
クラス設計と補完的に働いている。一方、フレーム変換関数側は「1 フレームを
受け取り 1 フレームを返す」という暗黙の契約に従う必要があり、
契約を型として明示できない Python では、関数の仕様をドキュメントで
共有することが重要になる。
